%!TEX root = ./../../thesis-trier.tex

\label{def:rankword}
\label{def:defectword}
\label{def:exclword}
\label{def:duplword}
\label{def:transitionmtx} 
\label{def:fwordu}
\IN{Boolean transition matrix!automaton}
\IN{duplicate set!of a word}
\IN{permutational letter/word}
\IN{simple idempotent letter/word} 
\IN{idempotent letter/word}
\IN{rank!of a word}
\IN{defect!of a word}
\IN{excluded set!of a word}
\IN{function induced by word}

Let $\semiautomat{A}{Q}{\Sigma}{\delta}$. Each word $u \in \Sigma^*$  \emph{induces a transformation} 
$\fword{\delta}{u} \colon Q \to Q$ by $\applyfunc{q}{\fword{\delta}{u}} = \applytrans{\delta}{q}{u}$ for all $q \in Q$. 
For $u, v \in \Sigma^*$ we have $\fword{\delta}{uv} = \concat{\fword{\delta}{u}}{\fword{\delta}{v}}$
and $\fword{\delta}{{uv}}^{-1} = \concat{\fword{\delta}{v}^{-1}}{\fword{\delta}{u}^{-1}}$.
Compare the latter equation with Equation~\eqref{prelim:eqn:delta_inverse}
wheres $u$ and $v$ are ``swapped'' in the expression on the right hand side.
This follows as in the $\delta$-notation the left-most word 
is applied first whereas in the composition of two functions, by our definition,
the right-most function is applied first.

We call a letter $x \in \Sigma$ (resp.~a word $u \in \Sigma^*$)
a \emph{permutational / idempotent / simple idempotent} letter (resp.~word) 
if $\fword{\delta}{x}$ (resp.~$\fword{\delta}{u}$) is a permutation/idempotent/simple idempotent.
The \emph{rank} (resp.~\emph{defect}, \emph{excluded set}, \emph{duplicate set}) of a word $u \in \Sigma^*$, denoted by $\rank{u}$ (resp.~$\defect{u}$, $\excl{u}$, $\dupl{u}$) 
is the rank (resp.~defect, excluded set, duplicate set) of~$\fword{\delta}{u}$. 
Furthermore, when using terminology from functions, transformations, and permutations, as cyclically permuting the states, the image of a set, transition matrix,
being a cycle, having a fixed point etc., in the context of a word (or letter) $u \in \Sigma^*$,
we refer to the induced transformation~$\fword{\delta}{u}$.
We also say that~$\fword{\delta}{u}$ is the \emph{action of $u$} on the states
or that~$u$ \emph{acts as $\fword{\delta}{u}$}. The same terminology transfers when speaking about properties of $\fword{\delta}{u}$,
for example, we say that~$u$ acts as a cycle if~$\fword{\delta}{u}$ is a cycle.

For $u \in \Sigma^*$, we write $[u]_{\mathcal A}$ (or just $[u]$ when there is no ambiguity to which automaton we refer)
for the transition matrix of $\fword{\delta}{u}$.
\Cref{prelim:fig:transitionmtxCerny} depicts transition matrices for the \Cerny automata $\mathcal C_n$
from \cref{intro:sec:cerny} w.r.t.~the natural order of its state set~$[n]$.
When  the specific ordering of the states is not important, we
assume some order to be taken without explicitly mentioning it.\footnote{The maps $u \mapsto \fword{\delta}{u}$
	and $u \mapsto [u]$ are morphisms from $\Sigma^*$ into $\mathcal T_Q$ and the set of $|Q| \times |Q|$ matrices, respectively.
	Thus, by \cref{prelim:prop:UMP}, they
	are uniquely determined by their images on the letters~$\Sigma$.}

\begin{figure}[th]
	\centering
	\[
	[b]_{\mathcal C_n} =
	\left( {\begin{array}{ccccc}
			0 & 1 &  0 & \cdots & 0\\
			0 & 0 &  1 & \cdots & 0\\
			\vdots & \vdots & \vdots & \ddots & \vdots\\
			0 & 0 & 0 & \cdots & 1 \\
			1 & 0 & 0 & \cdots & 0
	\end{array} } \right)
	\qquad 
	[a]_{\mathcal C_n} =
	\left( {\begin{array}{cccccc}
			0 & 1 &  0 & \cdots & 0 & 0 \\
			0 & 1 &  0 & \cdots & 0 & 0 \\
			\vdots & \vdots & \vdots & \ddots & \vdots\\
			0 & 0 & 0 & \cdots & 1 & 0 \\
			0 & 0 & 0 & \cdots & 0 & 1
	\end{array} } \right)
	\]
	\caption[The transition matrices for the \Cerny automata $\mathcal C_n$]{
		The transition matrices for the \Cerny automata $\mathcal C_n$
		from \cref{intro:sec:cerny} (\cref{intro:fig:cerny})
		where the states are ordered w.r.t.~the natural order on $[n]$.
	}
	\label{prelim:fig:transitionmtxCerny}
\end{figure} 

\label{def:transformation_semigroup}
The \emph{transformation semigroup of $\mathcal A$} is 
\[ \TS{\mathcal A} = \{ \fword{\delta}{u} \mid u \in \Sigma^+ \}
\] 
and the \emph{transformation monoid} is $\TM{\mathcal A} =  \{ \fword{\delta}{u} \mid u \in \Sigma^* \}$. 
\label{def:transformation_monoid}
\index{transformation semigroup!of automaton}
\index{transformation monoid!of automaton}

\begin{proposition}
	\label{sync_max:prop:grp_of_units_inv}	
	Let $\semiautomat{A}{Q}{\Sigma}{\delta}$
	and $\fword{\delta}{u}$ be a permutation for $u \in \Sigma^*$.
	Then $\fword{\delta}{u}^{-1} \in \TM{\mathcal A}$,
	i.e., the inverse permutation
	is induced by some word over $\Sigma^*$.
\end{proposition}
\begin{proof}
	By finiteness, 
    we have $\fword{\delta}{u^i} = \operatorname{id}_Q$ for some $i > 0$.
	Then $\fword{\delta}{u}^{-1} = \fword{\delta}{u^{i-1}} \in \TM{\mathcal A}$.	
\end{proof}

Next, we determine the group of units in $\TM{\mathcal A}$.

\begin{proposition}
	\label{sync_max:prop:grp_of_units_TA}	
	Let $\semiautomat{A}{Q}{\Sigma}{\delta}$ and 
	$\Gamma = \{ x \in \Sigma \mid x \text{ is a permutational letter } \}$.
	Then for the group of units $G$ of $\TM{\mathcal A}$ we have
	\[
	G = \{ \delta_u \mid \text{$\delta_u$ is a permutation and $u \in \Sigma^*$} \}
	= \{ \fword{\delta}{u} \mid u \in \Gamma^* \},
	\]
	and if~$G$ is non-trivial, equivalently $\Gamma \ne \emptyset$,
	we also have $G =  \{ \fword{\delta}{u} \mid u \in \Gamma^+ \}$.
\end{proposition}
\begin{proof}
	We have	$G \subseteq \{ \delta_u \mid \text{$\delta_u$ is a permutation and $u \in \Sigma^*$} \}$ 
	as every transformation in $G$, having an inverse, is a permutation.
	Next, by~\crefitem{prelim:prop:rank-defect}{prelim:prop:item:rank-eqn} 
	the composition of a non-permutation with a transformation
	gives a non-permutation.  So, for $u \in \Sigma^*$,
	$\fword{\delta}{u}$ is a permutation if and only if $u \in \Gamma^*$, i.e.,
	\[ 
	\{ \delta_u \mid \text{$\delta_u$ is a permutation and $u \in \Sigma^*$} \}
	= \{ \fword{\delta}{u} \mid u \in \Gamma^* \}.
	\]	
	By \cref{sync_max:prop:grp_of_units_inv},
	for $u \in \Gamma^*$
	we have $\fword{\delta}{u}^{-1} = \fword{\delta}{v}$ for some $v \in \Sigma^*$, i.e.,
	$\fword{\delta}{u}$ has an inverse in $\TM{\mathcal A}$
	and so $\fword{\delta}{u} \in G$.
	This gives $\{ \fword{\delta}{u} \mid u \in \Gamma^* \} \subseteq G$.
	Combining everything,
	we find $G = \{ \delta_u \mid \text{$\delta_u$ is a permutation and $u \in \Sigma^*$} \}
	= \{ \fword{\delta}{u} \mid u \in \Gamma^* \}$.
	
	Lastly, if $\Gamma \ne \emptyset$,
	then $\fword{\delta}{u}$ equals the identity transformation
	for some $u \in \Gamma^+$.
	This gives $\{ \fword{\delta}{u} \mid u \in \Gamma^* \}
	= \{ \fword{\delta}{u} \mid u \in \Gamma^+ \}$.
\end{proof}

Thus the group of units of $\TM{\mathcal A}$ is precisely the set of permutations in~$\TM{\mathcal A}$
and every permutation group in $\TS{\mathcal A}$ is contained in the group of units of $\TS{\mathcal A}$ (but not every
subset $S \subseteq \TM{\mathcal A}$ that forms a group in itself; for example
for every non-permutational word $u$ permuting its image $\delta(Q, u)$ 
the set $\{ \fword{\delta}{u^i} \mid i \ge 1 \}$
forms a group, but contains no permutation).